\documentclass[../../../thesis.tex]{subfiles}
\begin{document}
\subsection{Použitie umelej inteligencie}\label{sec:utilizingAI}
Na optimalizáciu formulácie myšlienok môžeme využiť služby
umelej inteligencie (AI, z~ang. \emph{\foreignlanguage{english}{artificial intelligence}}) a~tzv. veľkých jazykových modelov (LLM, z~ang. \emph{\foreignlanguage{english}{large language model}}).
Umelá inteligencia dokáže kontроlovať rozsiahlejšie časti prác,
vyhľadáva chyby a navrhuje vhodnejšie
formulácie na základe pravidiel,
ktoré sme aplikovali v predchádzajúcom texte.
Neosvedčuje sa však pri kompozícii textov.
Neuspokojivé výsledky dosahujeme aj v prípadoch,
kedy necháme umelú inteligenciu preformulovať celé odseky.
Zanáša do nich chyby a nezmysly, ktoré tam pôvodne neboli.
Ťažko sa potom odhaľujú.
Tento jav poznáme ako tzv.
halucinácie a trpia nimi všetky nástroje AI,
vrátane najznámejšieho ChatGPT.

Napriek tomu predstavujú služby AI silný nástroj pri tvorbe pôvodného obsahu, zvlášť užitočné sú tzv. generatívne umelé inteligencie (GAI),
ktoré dokáže vytvárať výstupy takmer na nerozoznanie od tvorby človeka.
Ich správna aplikácia nepochybne prispieva k~vyššej jazykovej a obsahovej kvalite záverečných prác.
Treba však mať na pamäti, že záverečná práca má byť
pôvodné autorské dielo študenta a všetky časti,
ktoré nepochádzajú od autora musia byť riadne
zdokumentované a deklarované v~zozname použitých zdrojov.
V žiadnom prípade sa neodporúča, aby
GAI formulovala pôvodné myšlienky
alebo súvislé časti práce.
Takéto konanie považujeme za nečestné podobne,
ako keby prácu písal niekto iný,
prípadne by boli celé odseky prebrané z iného zdroja bez korektného citovania (pozri kapitolu \ref{sec:citation}).

Používane umelej inteligencie pri písaní záverečných prác
upravuje opatrenie rektora STU v Bratislave č. 1/2024-O, ktoré budeme ďalej v~texte uvádzať ako \uv{opatrenie}~\cite{opatrenie12024}.

\subsubsection*{Povolené činnosti umelej inteligencie bez potreby deklarácie}
Podľa čl. V, ods. 2, písm. a) opatrenia môžu študenti používať GAI bez potreby deklarácie na tieto činnosti: kontrola gramatiky, oprava textu, tvorba osnovy, zhromažďovanie informácií a~použitie výpočtových metód a~softvérov, ktoré obsahujú prvky AI.

\subsubsection*{Deklarácia činnosti generatívnej umelej inteligencie}
Čl. V, ods. 2, písmeno b) opatrenia obsahuje zoznam možností použitia GAI, ktoré je potrebné v práci deklarovať na konci po zozname literatúry.
Ide o nasledujúce činnosti: preklady medzi jazykmi, úpravy a reformulácie textu, tvorba zhrnutia a rešerší, citovanie odpovedí GAI, tvorba počítačových programov, tvorba grafického obsahu a obrázkov.

V závere práce, uvedieme za zoznamom literatúry
časti textu vytvorené s~pomocou AI, spôsob ich využitia a použitý nástroj AI (Opatrenie \cite{opatrenie12024}, čl. VI., ods. 2).

V hlavnom súbore záverečnej práce \verb|thesis.tex| je príkaz
\verb|\FEIaiDeclaration| na načítanie súboru \verb|ai_declaration.tex| z~priečinka \verb|includes|.
Každý výskyt použitia nástrojov AI zapíšeme ako
položku \verb|\item| do pripraveného prostredia
\verb|trivlist| v tomto súbore.
Formát a~obsah jednotlivých záznamov je naznačený v prílohe opatrenia číslo 1/2024-O.
Záznamy obsahujú tieto prvky:
\begin{trivlist}
  \item Názov spoločnosti (dátum), Názov nástроja, čast práce, účel použitia.
\end{trivlist}

Predchádzajúci vzorec vygeneroval nástroj ChatGPT 4o od firmy
OpenAI dňa 2. 2. 2025 na základe analýzy spomínaného opatrenia.
V deklarácii použitia umelej inteligencie sa zapíšeme tento
záznam:
\begin{trivlist}
  \item OpenAI (2025), ChatGPT 4o, časť \ref{sec:utilizingAI},
  generovanie vzorca záznamu použitia AI.
\end{trivlist}

Súčasná verzia šablóny FEIstyle nedisponuje nástrojmi na
automatizáciu záznamov činnosti AI.
Preto ich treba zapisovať ručne do súboru
\verb|includes/ai_declaration.tex|.
\end{document}